\chapter{1905: Adolf von Baeyer}
\label{chap:chemistry1905}

The Nobel Prize in Chemistry for the year 1905 was awarded to Adolf von Baeyer.

\textbf{Motivation:}
in recognition of his services in the advancement of organic chemistry and the chemical industry, through his work on organic dyes and hydroaromatic compounds

\section{Adolf von Baeyer}

\subsection*{Early Life and Education}

Adolf von Baeyer was born on 1835-10-31 in Berlin, Prussia (now Germany).
Johann Friedrich Wilhelm Adolf von Baeyer was a German chemist who synthesized indigo and developed a nomenclature for cyclic compounds. He was ennobled in the Kingdom of Bavaria in 1885 and was the 1905 recipient of the Nobel Prize in Chemistry.

Baeyer studied under Robert Bunsen at Heidelberg and worked with August Kekulé in Heidelberg, Ghent, and Berlin, where he absorbed the structural theories that were transforming organic chemistry in the 1860s. After taking his doctorate in Berlin in 1858, he continued research with Kekulé and began the investigations of uric acid that established his reputation. In 1872 he was called to the University of Strasbourg, and in 1875, after the death of Justus von Liebig, he accepted the chair of chemistry at the University of Munich, which he occupied until his retirement in 1915.

\subsection*{Scientific Context}

The synthesis of natural coloring matters was one of the great ambitions of nineteenth-century organic chemistry. The commercial success of mauveine and the synthetic dyes that followed it had created a dye industry of enormous value, and the possibility of reproducing the colors of nature in the laboratory exercised the best chemists of the age. Indigo, the blue dye extracted from the indigo plant, was among the most prized, and its structure had resisted analysis since the beginnings of organic chemistry. Baeyer's success in this enterprise, achieved after years of effort, established a model for the way the most difficult problems of structure and synthesis could be solved.

\subsection*{Career and Major Research}

Baeyer was professor at Munich from 1875 until his retirement. He achieved important laboratory syntheses of indigo beginning in 1880, providing a foundation for its later industrial manufacture, and he established important structural relationships among uric acid derivatives. He developed the strain theory of carbon ring compounds and contributed to the chemistry of hydroaromatic compounds, work of great importance to the dye industry.

Baeyer's earliest important research, carried out under Kekulé's influence, was on uric acid, a constituent of urine whose derivatives later became central to purine chemistry. He clarified the relationships among uric acid and related compounds, work that helped prepare the way for later structural studies of purines. Together with his later synthesis of indigo, this research established him as one of the leading organic chemists of Europe before he was forty.

\subsection*{Nobel Prize-winning Work}

The work for which Adolf von Baeyer was awarded the Nobel Prize in Chemistry in 1905 was described by the Nobel Committee as: "in recognition of his services in the advancement of organic chemistry and the chemical industry, through his work on organic dyes and hydroaromatic compounds".

This research fundamentally advanced the field by synthesizing indigo and advancing the chemistry of cyclic compounds, including the development of systematic nomenclature for ring structures.

The synthesis of indigo was Baeyer's most celebrated achievement. He determined the structure of indigo by careful degradation experiments and then set about the far more difficult task of building the molecule from simple materials. In 1880 he announced an important synthesis of indigo from cinnamic acid, and in 1883 he published a structural formula for the dye. His laboratory syntheses and structural work provided the basis for later industrial processes, which were developed and commercialized by firms including the Badische Anilin- und Sodafabrik. Industrial synthetic indigo eventually displaced much of the centuries-old cultivation of the natural dye and transformed the dye trade.

\subsection*{The Work in Detail}

Baeyer's investigation of indigo was remarkable both for its analytical rigor and for its synthetic ambition. He degraded the natural dye into simpler substances, relating these to known compounds until he had fixed the position of every atom in the molecule, and he then checked his conclusions by total synthesis. His structural formula accounted for all the known transformations of indigo and predicted new ones, and it served as the working basis for the industrial processes that followed.

In the chemistry of cyclic compounds, Baeyer made two contributions of enduring importance. The first was his strain theory, published in 1885, which explained the stability of ring compounds by the deviation of their bond angles from the tetrahedral value. Small rings, he argued, are strained and hence reactive, while larger rings are stable. The second was his development of a nomenclature for cyclic compounds, including the use of the prefixes cycl- and the system of numbering the atoms of a ring, conventions still in use today. His later work on hydroaromatic compounds, the hydrogenated derivatives of aromatic rings, grew out of his interest in the terpenes and carried his ideas into a new field.

\subsection*{Reception}

Baeyer's achievements were recognized in Germany and abroad with every honor the scientific world could bestow. He was ennobled by the Kingdom of Bavaria in 1885, becoming Adolf von Baeyer, and he received honorary degrees and academy memberships throughout Europe. His Munich laboratory became the most famous school of organic chemistry in the world, and it was said that no other laboratory of the time produced so many distinguished chemists. The award of the Nobel Prize in 1905, at the age of seventy, was the culmination of a career that had already made him a legendary figure.

\subsection*{Significance and Impact}

The impact of Adolf von Baeyer's work extends far beyond the specific achievement recognized by the Nobel Prize.
Adolf von Baeyer's synthesis of indigo revolutionized the dye industry, and Adolf von Baeyer's systematic work on cyclic compounds established principles of ring strain and stereochemistry that remain fundamental.

The industrial production of synthetic indigo, begun in 1897, quickly displaced the natural product and destroyed an agricultural industry that had flourished for centuries in India and elsewhere. The same pattern was repeated for other natural dyes and, more broadly, for the whole range of synthetic organic products that the German chemical industry came to supply. Baeyer's methods and his approach to structure determination set the standard for the synthesis of natural products, and his strain theory provided a rationale for the reactivity of small-ring compounds that remains valid.

\subsection*{Later Career and Honors}

Adolf von Baeyer continued his scientific work until 1917-08-20, passing away in Starnberg, Germany.
He received numerous honors including membership in major scientific academies and societies.
Adolf von Baeyer was ennobled by the Kingdom of Bavaria and served as president of the German Chemical Society.

In addition to the Nobel Prize, Baeyer received the Davy Medal of the Royal Society in 1881 and the Faraday Lectureship in 1907, and he was elected to the Bavarian Academy of Sciences, of which he served as president. His later years were darkened by the First World War, and he died in 1917 in Starnberg, near Munich. The laboratory he built and the school of chemists he trained had made Munich the undisputed capital of organic chemistry in Europe.

\subsection*{Legacy}

The legacy of Adolf von Baeyer endures through the lasting impact of his scientific contributions.
Adolf von Baeyer's strain theory for ring compounds and nomenclature systems for cyclic structures remain essential in organic chemistry. The von Baeyer name is associated with the Baeyer-Villiger oxidation, a key synthetic transformation.

Baeyer's influence is still visible in the daily practice of organic chemistry. The concepts of ring strain and of angle strain are used to explain the reactivity of cycloalkanes and of strained polycyclic molecules, and his nomenclature for ring compounds is part of the international system of chemical nomenclature. The Baeyer strain theory, modified by the later recognition that rings are not flat, remains the starting point for understanding the conformations of cyclic molecules. His name is also attached to the Baeyer-Villiger oxidation, a reaction discovered by his student Victor Villiger, which is still widely used in synthesis.

\subsection*{Collaborators and Students}

Baeyer's students formed an extraordinary group of later leaders of chemistry. They included Emil Fischer, the 1902 Nobel laureate whose work on sugars and purines carried on Baeyer's researches on uric acid; Richard Willstätter, the 1915 Nobel laureate; and Victor Villiger, Baeyer's assistant, with whom he discovered the oxidation reaction that bears their names. The standard of technique and the insistence on rigorous proof of structure that Baeyer demanded in his laboratory were transmitted by his students to laboratories throughout the world.

\subsection*{Modern Developments}

Baeyer's synthetic chemistry, including his synthesis of indigo, laid the groundwork for the modern dye and pigment industry, which still produces billions of dollars of colorants annually. The Baeyer strain theory and his structural investigations informed the development of cyclic molecules, including modern pharmaceuticals and agrochemicals. The Baeyer–Villiger oxidation, named after his student, remains a standard synthetic tool.

The chemistry of small-ring compounds, which Baeyer opened with his strain theory, has become a rich field of modern research: cyclopropane rings appear in natural products and pharmaceuticals, strained molecules such as cubane are studied as novel materials and energetic compounds, and the concepts he introduced are used to design molecules with predictable reactivity. The dye industry that he helped to create now produces colors for textiles, printing, electronics, and biological imaging, and the structural methods he pioneered remain central to the identification and synthesis of new compounds.

\subsection*{Selected Publications}

"Synthese des Indigos" (1880)
"Über die Constitution des Indigos" (1883)
"Chemie der cyclischen Verbindungen" (1885)
